\documentclass[11pt]{article}
\usepackage[margin=1in]{geometry}
\usepackage{amsmath,amssymb,amsthm}
\usepackage{graphicx}
\usepackage[hidelinks]{hyperref}

\newtheorem*{theorem}{Theorem}
\newtheorem*{conclusion}{Conclusion}

\newcommand{\Hh}{\mathcal H}
\newcommand{\N}{\mathbb N}
\newcommand{\ip}[2]{\langle #1,#2\rangle}

\title{Complete sequences are weighted lower semi-frames}
\author{Literature-implied status packet for arXiv:2606.00764}
\date{June 8, 2026}

\begin{document}
\maketitle

\section*{Open question}

The source paper is E. Zikkos, \emph{The Gaussian Gabor system at the
critical density is a weighted lower semi frame}, arXiv:2606.00764.
It quotes Conjecture 3.1 of P. Balazs, M. Corso, and D. T. Stoeva,
\emph{Weighted frames, weighted lower semi frames and unconditionally
convergent multipliers}, arXiv:2310.18957.

The conjecture asks whether every complete family in a separable
Hilbert space is a weighted lower semi-frame. The source paper proves
the case where the family contains a complete minimal subfamily, then
leaves open the case in which every complete subfamily has infinite
excess. It also singles out the Muntz family
\(\{t^{\lambda_n}\}\subset L^2(0,1)\), with
\(\sum_n \lambda_n^{-1}=\infty\), as an unknown example.

\begin{center}
\includegraphics[width=.92\linewidth]{figures/open_problem_crop.png}
\end{center}

\section*{Literature-implied answer}

N. Ozawa's preprint \emph{Total sequences and the lower frame
inequality}, arXiv:2507.04899, proves exactly the required Hilbert-space
statement. Ozawa states the theorem as a solution to a MathOverflow
exercise, not as a solution of the Balazs--Corso--Stoeva or Zikkos
formulation; the match below is the reformulation recorded by this
packet.

\begin{theorem}[Ozawa, arXiv:2507.04899]
Let \((v_n)_{n=1}^\infty\) be a total sequence in a Hilbert space \(H\).
Then there is a sequence \((\lambda_n)_{n=1}^\infty\) of positive
numbers such that
\[
  \|x\|^2 \leq \sum_{n=1}^\infty \lambda_n |\ip{x}{v_n}|^2
\]
for every \(x\in H\), with the right-hand side possibly equal to
\(\infty\).
\end{theorem}

This is precisely the lower semi-frame inequality after scaling:
define
\[
  \omega_n=\sqrt{\lambda_n}.
\]
Then
\[
  \sum_{n=1}^\infty |\ip{x}{\omega_n v_n}|^2
  =
  \sum_{n=1}^\infty \lambda_n |\ip{x}{v_n}|^2
  \geq \|x\|^2
\]
for every \(x\in H\). Thus \((\omega_n v_n)\) is a lower semi-frame
with lower bound \(1\). If the BCS formulation requests an arbitrary
prescribed lower bound \(A>0\), replace \(\lambda_n\) by \(A\lambda_n\).

\begin{conclusion}
Ozawa's theorem fully answers the conjecture quoted in arXiv:2606.00764
affirmatively. In particular, the Muntz-family example in the source
paper is weighted lower semi-frame whenever it is complete.
\end{conclusion}

\section*{Proof idea from Ozawa}

Let
\[
  H_n=\overline{\operatorname{span}}\{v_k:k\geq n\},
  \qquad
  H_\infty=\bigcap_n H_n.
\]
The descending chain \(H_1\supset H_2\supset\cdots\) has one-dimensional
successive quotients at most:
\[
  \dim(H_n\cap H_{n+1}^{\perp})\leq 1.
\]
Choose an orthonormal basis from these quotient directions together
with an orthonormal basis of \(H_\infty\), ordered so that the \(n\)-th
basis vector lies far enough in the tail space \(H_{\lceil n/2\rceil}\).
Approximate those basis vectors by finite linear combinations of the
tail vectors \(v_k\). The approximants define an operator that is a
small-norm perturbation of the identity, hence invertible. Expanding the
resulting strong-operator identity back into the original \(v_k\)'s and
applying Cauchy--Schwarz produces the coefficients \(\lambda_k\).

\section*{Status}

This is a literature-implied answer packet, not an original pipeline
solution. The earlier partial packet proving the countable complete-block
criterion remains valid, but it is superseded for the full conjecture by
arXiv:2507.04899.

\end{document}
